\section{Summary}
\label{sec:summary}

\paragraph{What this report is for.} It gives the thermal profile, the drag and the air loads of the two
candidate 720\,mm nose cones. It compares design A with design B, and gives the structure team (ST)
what it needs to size the shell, the aluminium tip insert and the joint to the body tube. ST computes the
stresses from these loads.

\paragraph{The two noses.} Both are 720\,mm long and \SI{180}{mm} in diameter, with a \SI{4.5}{mm}
rounded tip. Their geometry is in \file{2027_C_FD_NOSE-CONE-720MM_GEOMETRY}.

\begin{table}[H]
\centering
\small
\begin{tabular}{@{}lll@{}}
\toprule
 & Design A (baseline) & Design B (backup)\\
\midrule
Shape & LV-Haack & Von K\'arm\'an\\
3U payload & inside the nose & elsewhere in the rocket\\
Curved part + cylindrical coupler & \val{A-xS} + \val{A-c}\,mm & \val{B-xS} + \val{B-c}\,mm\\
\bottomrule
\end{tabular}
\end{table}

\begin{keybox}[Main findings]
\begin{enumerate}[leftmargin=*, itemsep=4pt]
\item \textbf{Thermal profile (Section~\ref{sec:heating}).} The shell is hottest just behind the tip and
cools along the nose. Its \textbf{outer face} stays below the \SI{120}{\celsius} of the DP490 bond: at
most \val{th-B-nominal-Tout-hi}\,\si{\celsius} on a nominal day and \val{th-B-hot-Tout-hi}\,\si{\celsius}
on a hot day, so the hot-day margin is small. Its \textbf{inner face} is above the \SI{67}{\celsius} of
Req\_10 over the whole length, already on a nominal day (\val{th-A-nominal-Tin-lo} to
\val{th-B-nominal-Tin-hi}\,\si{\celsius}); so is the air inside.
\item \textbf{A against B, heating.} Both designs heat up the same way: the shell temperatures differ by
less than \SI{1}{\celsius}. Only the tip insert of B is 5--6\,\si{\celsius} cooler.
\item \textbf{A against B, drag (Section~\ref{sec:drag}).} Below Mach 1 both noses have the same drag.
Above Mach 1 design B has \val{BlessA-M120}--\val{BlessA-M147}\,\% less: \val{B-F-M147}\,N against
\val{A-F-M147}\,N at burnout (Fluent). B reaches \val{B-apogee}\,m against \val{A-apogee}\,m for A.
\item \textbf{The largest structural load comes from the nose being sealed.} The air inside stays at the
pad pressure (\val{pad-p}\,kPa), while the outside pressure falls to \val{A-nominal-LC3-pinf}\,kPa at
apogee. The shell is then pushed outward by about \absval{A-nominal-LC3-dpmin}\,kPa, and the nose is pulled
off the joint with up to \absval{B-hot-LC3-Fphi}\,N, close to the 2000\,N of Req\_07. Vent holes would remove
this load.
\item \textbf{The highest inward pressure is on the tip at burnout}: up to \val{B-hot-LC1-dpstag}\,kPa on
the rounded tip only. Further back it is at most \val{A-hot-LC1-dpmax2}\,kPa (A) and
\val{B-hot-LC1-dpmax2}\,kPa (B), and about 10\,kPa or less beyond 50\,mm.
\item \textbf{Side force}: \val{B-hot-LC1-Nperdeg}\,N per degree of angle of attack at burnout (hot day),
acting \val{A-xcp}\,mm (A) or \val{B-xcp}\,mm (B) behind the tip.
\end{enumerate}
\end{keybox}

\begin{table}[H]
\centering
\caption{Values to design for. Each value is the worst case over the flight. The hot day is the worst
case for every quantity. These are flight values, without a factor of safety, and air loads only: the
inertial loads and the recovery load (Req\_07) are not included. \ok: below the limit; \bad: above it.}
\label{tab:design}
\small
\begin{tabular}{@{}llcclc@{}}
\toprule
 & Unit & Design A & Design B & When & Limit\\
\midrule
\grouprow{6}{Temperatures (Section~\ref{sec:heating})}
Shell outer face, 10\,mm (hottest station) & \si{\celsius} & \val{th-A-hot-Tout-hi} & \val{th-B-hot-Tout-hi} & 9.4\,s & 120 (DP490) \ok\\
Shell outer face, shoulder & \si{\celsius} & \val{th-A-hot-Tout-shoulder} & \val{th-B-hot-Tout-shoulder} & 10.2\,s & 120 (DP490) \ok\\
Shell inner face, 10\,mm (hottest station) & \si{\celsius} & \val{th-A-hot-Tin-hi} & \val{th-B-hot-Tin-hi} & 11.2\,s & 67 (Req\_10) \bad\\
Shell inner face, shoulder & \si{\celsius} & \val{th-A-hot-Tin-shoulder} & \val{th-B-hot-Tin-shoulder} & 12.2\,s & 67 (Req\_10) \bad\\
Air inside the nose (upper bound) & \si{\celsius} & \val{A-hot-Tair} & \val{B-hot-Tair} & 12\,s & 67 (Req\_10) \bad\\
Aluminium tip insert & \si{\celsius} & \val{A-hot-Ttip} & \val{B-hot-Ttip} & 12--13\,s & \\
\midrule
\grouprow{6}{Drag and forces on the joint (Sections~\ref{sec:drag} and~\ref{sec:loads})}
Drag (compression) & N & \val{A-hot-LC1-Dn} & \val{B-hot-LC1-Dn} & burnout & \\
Pull-off force of the sealed nose (tension) & N & \absval{A-hot-LC3-Fphi} & \absval{B-hot-LC3-Fphi} & apogee & Req\_07: 2000\\
Side force per degree of angle of attack & N/$^\circ$ & \val{A-hot-LC1-Nperdeg} & \val{B-hot-LC1-Nperdeg} & burnout & \\
Where the side force acts, from the tip & mm & \val{A-xcp} & \val{B-xcp} & always & \\
\midrule
\grouprow{6}{Pressure on the shell (Section~\ref{sec:loads})}
Inward pressure on the rounded tip & kPa & \val{A-hot-LC1-dpstag} & \val{B-hot-LC1-dpstag} & burnout & \\
Inward pressure behind the tip ($x > 2$\,mm) & kPa & \val{A-hot-LC1-dpmax2} & \val{B-hot-LC1-dpmax2} & burnout & \\
Outward pressure on the shell & kPa & \absval{A-hot-dpminall} & \absval{B-hot-dpminall} & coast to apogee & \\
\bottomrule
\end{tabular}
\end{table}

\paragraph{Where the numbers come from.} Eight Fluent CFD runs of the two noses (2D, no angle of attack)
at three moments of the climb, and the team's Python flight model, calibrated on these runs, for the
rest of the flight and for the hot day. Appendix~\ref{app:method} explains the method. The data behind
every figure are in the CSV files of \file{report_720/data/}.
